\section{Discussion and conclusions}
\label{sec:conclusions}

A reservoir requirement is precise only once we specify what must be
reproduced. The results separate three ingredients that are often
merged: the reservoir capacity, its calibration, and the observable
being tested. For the complete canonical law at coexistence, capacity
and calibration interact: the secant calibration removes the phase bias
caused by the latent heat, so the capacity needs to control only the
tilt within each phase. This is why two-phase coexistence needs $c_N\gg
N^{3/2}$ rather than the $N^2$ suggested by the latent heat or by the
subsystem's canonical heat capacity. Three distinct phase energies,
including the combined energy of two coexisting copies, restore the
$N^2$ requirement. Smooth reservoirs with positive heat capacity obey
the same curvature mechanism when a suitable calibration and tail
control are available. The heat-capacity convention also matters at
finite size, as finite ideal-gas contact calculations show
\cite{CortiEtAl2023}.

The observable matters as much. Phase weights should be reported with
any optimized error, since at the threshold scale the smallest
full-state TV can come from discarding a phase. Marginal and joint
accuracy differ: with balanced phase weights, a shared reservoir can
leave each copy canonical while anticorrelating their phases, an
equilibrium effect that needs no direct interaction between the copies.
Other criteria are not equivalent to full-state accuracy: a small
Euclidean histogram norm or macroscopic energy concentration does not
imply it, and the log-density barrier between a phase center and the
interphase midpoint can remain lowered by a large amount while the
full-state TV tends to zero.

The proved scope is equilibrium coupling to the stated reservoir with
additive energies, fixed positive limiting phase weights, and the
explicit fluctuation and tail hypotheses of each theorem. The
short-range results require a fixed number of colors $q\ge q_0(d)$.
Their threshold-scale law holds for every fixed capacity coefficient
$\gamma>0$ in dimensions above two, and for sufficiently large $\gamma$
in two dimensions. The square-torus interfacial calculation is a stated
capillarity model, not a derivation of microscopic droplet or slab
morphology, and the Gaussian, capillarity, and linear-capacity
appendices analyze stated models exactly, whereas the contour appendix
proves a microscopic hypothesis.

These limits suggest open questions: the two-dimensional threshold-scale
law for small $\gamma$, where interfacial configurations compete with
the phases; Potts models with fewer colors and other interacting
systems; interacting spatial subregions of a single system, which a
shared-reservoir model of independent copies does not describe; quantum
systems, in particular observables that do not commute with the energy;
and explicit finite-size error bounds.
